\documentclass[11pt]{article}
\usepackage[margin=1in]{geometry}
\usepackage{amsmath,amssymb}
\usepackage{enumitem}
\usepackage{bookmark}
\hypersetup{
  pdftitle={PSET 5: Functions of several variables and optimization with several variables},
  hidelinks}

\title{PSET 5: Functions of several variables and optimization with
several variables}
\author{}
\date{\vspace{-2.5em}}
 
\newcounter{problem}
\newcommand{\problem}[1]{\stepcounter{problem}\section*{Problem \theproblem: #1}}

\begin{document}
\maketitle

% Shared assignment questions (header block).
% Include at the top of any problem set with:  % Shared assignment questions (header block).
% Include at the top of any problem set with:  % Shared assignment questions (header block).
% Include at the top of any problem set with:  \input{assignment-qs}
\section*{Assignment Qs}

\begin{itemize}
\item Name
\item How long did this problem set take you (round to nearest hour)? \quad
  0-1 hour \ 2-3 hours \ 4-5 hours \ 5+ hours
\item How difficult was this problem set? \quad
  very easy \ 1 \ 2 \ 3 \ 4 \ 5 \ very challenging
\end{itemize}

\section*{Assignment Qs}

\begin{itemize}
\item Name
\item How long did this problem set take you (round to nearest hour)? \quad
  0-1 hour \ 2-3 hours \ 4-5 hours \ 5+ hours
\item How difficult was this problem set? \quad
  very easy \ 1 \ 2 \ 3 \ 4 \ 5 \ very challenging
\end{itemize}

\section*{Assignment Qs}

\begin{itemize}
\item Name
\item How long did this problem set take you (round to nearest hour)? \quad
  0-1 hour \ 2-3 hours \ 4-5 hours \ 5+ hours
\item How difficult was this problem set? \quad
  very easy \ 1 \ 2 \ 3 \ 4 \ 5 \ very challenging
\end{itemize}


\problem{Find first partial derivatives}

Find all of the first partial derivatives of each
function.\footnote{Inspired by Grimmer HW6.3}

\begin{enumerate}[label=\alph*.]
\item $f(x,y) = \frac{1}{3}x^5 + 3x^3y^2 + 3xy^4$
\item $g(x,y) = 4xe^{4y+1}$
\item $k(x,y) = \dfrac{x-y}{x+y}$
\end{enumerate}

\problem{Find the gradient, Hessian and Jacobian}

Find the gradient, Hessian $H$, and Jacobian, $J$, for the following functions.\footnote{Inspired by
  Grimmer HW7.3} Evaluate the gradient at the given points.

\begin{enumerate}[label=\alph*.]
\item $g(x,y) = x^4+y^4-4xy+2$
  \begin{enumerate}[label=\roman*.]
  \item Gradient: (evaluate at: \quad $(x,y) = (1,1)$)
  \item Jacobian:
  \item Hessian:
  \end{enumerate}
\item $k(x,y) = (1+xy)(x+y)$
  \begin{enumerate}[label=\roman*.]
  \item Gradient: (evaluate at: \quad $(x,y) = (1,-1)$)
  \item Jacobian:
  \item Hessian:
  \end{enumerate}
\end{enumerate}

\problem{Find the critical points}

Find the local minimum values, local maximum values, and saddle point(s)
of the function. Remember the process we discussed in class: calculate
the gradient, set it equal to zero to solve the system of equations,
calculate the Hessian, and assess the Hessian at critical values. Be sure
to show your work on each of these steps.\footnote{Inspired by Grimmer
  HW7.4}

\begin{enumerate}[label=\alph*.]
\item $h(x,y) = x^4+y^4$
\item $g(x,y) = x^4+y^4-4xy+2$ (You can reuse your calculations from above)
\item $k(x,y) = (1+xy)(x+y)$ (You can reuse your calculations from above)
\end{enumerate}


\problem{Definite and Indefinite integrals}

The remaining problems cover integration. Solve the following integrals
using the antiderivative method.\footnote{Inspired by Gill 5.10, Grimmer
  HW4.1, Gill 5.13, and 5.14}

For all these problems, the basic approach to compute the definite
integral of $f(x)$ from $a$ to $b$ is by using the formula $F(b) - F(a)$,
where $F(x)$ is the \textbf{antiderivative} of $f$.

\begin{enumerate}[label=\alph*.]
\item $\int_{-1}^{0} (3x^2 - 3) \,dx$
\item $\int_2^4 e^{3y} \,dy$
\item $\int_3^3 \sqrt{x^5 - 2} \,dx$
\item $\int \left(x^2 - x^{-\frac{1}{2}}\right) \,dx$
\end{enumerate}


\problem{Substitution and integration by parts}
\begin{enumerate}[label=\alph*.]
\item $\int 2x(x^2+1)^4 \,dx$
\item $\int \dfrac{3x^2}{x^3+5} \,dx$
\item $\int x e^{2x} \,dx$

\item Without solving, state whether you would use substitution or
  integration by parts for each, and why:
  \begin{enumerate}[label=\roman*.]
  \item $\int 5x^4 \left(x^5 - 2\right)^{3} \,dx$
  \item $\int x^2 \log(x) \,dx$
  \end{enumerate}
\end{enumerate}

\problem{Improper and multiple integrals}

For (a)--(c), determine whether the integral is convergent or divergent,
and evaluate those that are convergent.\footnote{Inspired by Grimmer
  HW4.3} For (d)--(e), evaluate the iterated integral. In each case,
state whether the order of integration is forced, and why.

\begin{enumerate}[label=\alph*.]
\item $\int_1^{\infty} \left(\dfrac{1}{3x}\right)^2 \,dx$
\item $\int_0^{\infty} e^{-x} \,dx$
\item $\int_{-\infty}^0 x^3 \,dx$
\item $\int_{0}^1 \int_{2}^{3} x^2y^3 \,dx \,dy$
\item $\int_{0}^1 \int_0^{\sqrt{1-x^2}} 2x^3y \,dy \,dx$
\end{enumerate}

\problem{Interpretation}
No new computation is required for this problem. Answer each part in two to
four sentences of prose.

\begin{enumerate}[label=\alph*.]
\item \textbf{Reading a partial derivative.} Suppose we regress presidential
  approval in month $i$ on the employment rate and the price of gas, and
  obtain
  \[
    \text{Approval}_{i} = 0.8 + 0.5\,\text{Employ}_{i} - 0.25\,\text{Gas}_{i}.
  \]
  Write down $\dfrac{\partial\, \text{Approval}_{i}}{\partial\, \text{Employ}_{i}}$
  and interpret it substantively. Your interpretation should make clear what
  is happening to $\text{Gas}_{i}$ while employment changes, and why that
  matters for the claim you are making.

\item \textbf{Gradient versus Hessian.} Explain why the condition
  $\nabla f(\mathbf{a}) = \mathbf{0}$ is \emph{necessary} but not
  \emph{sufficient} for $\mathbf{a}$ to be a local extremum. Use your answers
  to Problem 3 to illustrate: both $(0,0)$ and $(1,1)$ have zero gradient,
  but they are different kinds of point. What does the Hessian tell you at
  each that the gradient could not? Finally, explain what it means when the
  determinant test returns $AC - B^2 = 0$, and why we cannot classify the
  point in that case.

\item \textbf{Integrals and probability.} A continuous probability density
  $f$ must satisfy $\int_{-\infty}^{\infty} f(x) \,dx = 1$. Explain what this
  requirement means, and what $\int_{a}^{b} f(x)\,dx$ represents for some
  interval $[a,b]$.
  \end{enumerate}

% Shared AI and Resources statement.
% Include in any problem set with:  % Shared AI and Resources statement.
% Include in any problem set with:  % Shared AI and Resources statement.
% Include in any problem set with:  \input{ai-statement}
\section*{AI and Resources statement}

Please list (in detail) all resources you used for this assignment. If
you worked with people, list them here as well. It is not enough to say
that you used a resource for help, you need to be specific on the link
and \emph{how} it was helpful. W/R/T gen AI tools (including GPT,
Claude, etc.) you cannot use them to do work on your behalf---you
cannot put in any of the questions, etc. You can ask for help on logic
/ sample problems. If you do use GPT or other AI tools, you need to
provide a link to your chat transcript. Any suspected academic
integrity violations will be immediately reported to the Dean of
Students.

\section*{AI and Resources statement}

Please list (in detail) all resources you used for this assignment. If
you worked with people, list them here as well. It is not enough to say
that you used a resource for help, you need to be specific on the link
and \emph{how} it was helpful. W/R/T gen AI tools (including GPT,
Claude, etc.) you cannot use them to do work on your behalf---you
cannot put in any of the questions, etc. You can ask for help on logic
/ sample problems. If you do use GPT or other AI tools, you need to
provide a link to your chat transcript. Any suspected academic
integrity violations will be immediately reported to the Dean of
Students.

\section*{AI and Resources statement}

Please list (in detail) all resources you used for this assignment. If
you worked with people, list them here as well. It is not enough to say
that you used a resource for help, you need to be specific on the link
and \emph{how} it was helpful. W/R/T gen AI tools (including GPT,
Claude, etc.) you cannot use them to do work on your behalf---you
cannot put in any of the questions, etc. You can ask for help on logic
/ sample problems. If you do use GPT or other AI tools, you need to
provide a link to your chat transcript. Any suspected academic
integrity violations will be immediately reported to the Dean of
Students.


\end{document}
